% 参考文献。FSS2026サンプルに合わせ thebibliography 環境で直接記述する。
% note=[要確認] の書誌は原論文PDFで最終確認すること（refs.bib 参照）。
% 並び順: 本文中で \cite が登場する順（01_intro → 02_related_work → 03_method）。
\begin{thebibliography}{99}

% 01_intro: 認知機能評価支援システム
\bibitem{kubota2018cognitive}
\colorbox{tbd}{久保田, 大保, 松田, 万治, 青村, 武田}:  % 修正: 名字のみに統一(黄色ハイライト)
自動車運転再開にむけた認知機能評価支援システム, 2018.

% 01_intro: 高齢者のためのスマートデバイス連動型シニアカー
\bibitem{nakamura2015seniorcar}
\colorbox{tbd}{中村, 田松, 久保田, 坂下}:  % 修正: 名字のみに統一(黄色ハイライト)
高齢者のためのスマートデバイス連動型シニアカー,
ロボティクス・メカトロニクス講演会 (ROBOMECH2015), 2A1-005, 2015.

% 01_intro: タッチ端末操作の操作特性の加齢変化
\bibitem{hatakenaka2024touch}
\colorbox{tbd}{畠中, 高橋, 下村, 原田}:  % 修正: 名字のみに統一(黄色ハイライト)
タッチ端末操作における操作特性の加齢変化 ―2024年度計測結果の報告―,
人間工学, Vol. 61, Suppl., 2024.

% 01_intro: 高齢者の認知特性とWebインタフェース
\bibitem{suzuki2009elderly}
\colorbox{tbd}{鈴木, 本宮, 原, 須藤, 佐藤, 熊田, 北島}:  % 修正: 名字のみに統一(黄色ハイライト)
高齢者の認知特性に適合した Web インタフェースのデザインに関する研究(2) ― 高齢者の認知機能が Web 操作に及ぼす影響 ―,
ヒューマンインタフェースシンポジウム 2009, 2009.

% 02_related_work: 高齢者にも優しい視線入力システム
\bibitem{murata2009gaze}
\colorbox{tbd}{村田, 早見, 森}:  % 修正: 名字のみに統一(黄色ハイライト)
視線入力システムの高性能化・実用化に関する研究 ― 高齢者にも優しい視線入力システムの開発 ―,
岡山大学大学院自然科学研究科, 2009.

% 02_related_work: VR環境での触感を伴ったフリック入力
\bibitem{kiyohara2022vrflick}
\colorbox{tbd}{清原, 沢田, 野呂}:  % 修正: 名字のみに統一(黄色ハイライト)
VR 環境における文字入力のための触感を伴ったフリック入力方法,
日本ソフトウェア科学会第39回大会 (2022年度) 講演論文集, 2022.

% 02_related_work: BubbleGlide（スマートウォッチ向け日本語曖昧入力） % 追加
\bibitem{toba2020bubbleglide}
\colorbox{tbd}{戸羽, 加藤, 山本}:  % 修正: 名字のみに統一(黄色ハイライト)
BubbleGlide: N-gram を用いたスマートウォッチ向け日本語曖昧入力ジェスチャーキーボード,
情報処理学会 インタラクション 2020 (IPSJ Interaction 2020), 2020.

% 02_related_work: 母音・子音の順に選択する間接タッチ用かな入力（WISS 2024） % 追加
\bibitem{wada2024vowelconsonant}
\colorbox{tbd}{和田, 白根, 崔, 志築}:  % 修正: 名字のみに統一(黄色ハイライト)
母音, 子音の順に選択を行う間接タッチ用かな文字入力手法,
インタラクティブシステムとソフトウェアに関するワークショップ (WISS 2024), 2024.

% 02_related_work: InvisiVowel
\bibitem{sugawara2023invisivowel}
\colorbox{tbd}{菅原, 鳴海}:  % 修正: 名字のみに統一(黄色ハイライト)
InvisiVowel: スマートフォンにおけるフリックの母音情報のみを利用したキーボード表示のない入力システム,
第31回インタラクティブシステムとソフトウェアに関するワークショップ (WISS 2023), 2023.

% 02_related_work: 母音列制約付き子音復元（VCCR）
\bibitem{hashimoto2026vccr}
\colorbox{tbd}{橋本}:  % 修正: 名字のみに統一(黄色ハイライト)
機械読唇の精度向上に向けた自然言語処理による母音列からの日本語文復元,
早稲田大学大学院基幹理工学研究科 修士論文, 2026.

% 02_related_work: KSPC（打鍵負担の指標） % 追加
\bibitem{mackenzie2002kspc}
\colorbox{tbd}{I. S. MacKenzie}:  % 追加: イニシャル+名字形式(黄色ハイライト)
KSPC (Keystrokes per Character) as a Characteristic of Text Entry Techniques,
Proceedings of the Fourth International Symposium on Human Computer Interaction with Mobile Devices (LNCS 2411), pp. 195--210, 2002.  % 追加

% 03_method: 学習データ CC-100 (XLM-R) % 追加
\bibitem{conneau2020xlmr}
\colorbox{tbd}{A. Conneau, K. Khandelwal, N. Goyal, V. Chaudhary, G. Wenzek, F. Guzmán, E. Grave, M. Ott, L. Zettlemoyer, V. Stoyanov}:  % 追加: イニシャル+名字形式
Unsupervised Cross-lingual Representation Learning at Scale,
Proceedings of the 58th Annual Meeting of the Association for Computational Linguistics (ACL), pp. 8440--8451, 2020.

% 03_method: 学習データ OpenSubtitles % 追加
\bibitem{lison2016opensubtitles}
\colorbox{tbd}{P. Lison, J. Tiedemann}:  % 追加: イニシャル+名字形式
OpenSubtitles2016: Extracting Large Parallel Corpora from Movie and TV Subtitles,
Proceedings of the Tenth International Conference on Language Resources and Evaluation (LREC 2016), pp. 923--929, 2016.

% 03_method: 学習データ JESC % 追加
\bibitem{pryzant2018jesc}
\colorbox{tbd}{R. Pryzant, Y. Chung, D. Jurafsky, D. Britz}:  % 追加: イニシャル+名字形式
JESC: Japanese-English Subtitle Corpus,
Proceedings of the Eleventh International Conference on Language Resources and Evaluation (LREC 2018), 2018.

% 03_method: LoRA
\bibitem{hu2022lora}
\colorbox{tbd}{E. Hu,} \colorbox{tbd}{Y. Shen,} \colorbox{tbd}{P. Wallis,} \colorbox{tbd}{Z. Allen-Zhu,} \colorbox{tbd}{Y. Li,} \colorbox{tbd}{S. Wang,} \colorbox{tbd}{L. Wang,} \colorbox{tbd}{W. Chen}:  % 修正: イニシャル+名字形式に統一(黄色ハイライト、各著者ごとに囲み行末で折り返す)
LoRA: Low-Rank Adaptation of Large Language Models,
International Conference on Learning Representations (ICLR), 2022.

\end{thebibliography}
